\documentclass[a4paper, 12pt]{article}

%%% SST LAB PROTOCOLL PREAMBLE
%%% 2019
%%%%%%%%%%%%%%%%%%%%%%%%%%%%%%%


%%% PACKAGES
%%%%%%%%%%%%%%%%%%%%%%%%%%%

\usepackage[ngerman]{babel}

\usepackage[utf8]{inputenc}
\usepackage{amsmath}
\usepackage{pgfplots}
\usepackage{tikz}
\usepackage[many]{tcolorbox}
\usepackage{graphicx}
\graphicspath{ {./graphics/} }
\usepackage{pdfpages}
\usepackage{dashrule}
\usepackage{float}
\usepackage{siunitx}
\usepackage{trfsigns}
\usepackage{booktabs}
\usepackage[european]{circuitikz}
\usepackage{tcolorbox}

%%% DOCUMENT GEOMETRY
%%%%%%%%%%%%%%%%%%%%%%%%%%%

\usepackage{geometry}
\geometry{
 a4paper,
 total={0.6180339887498948\paperwidth,0.6180339887498948\paperheight},
 top = 0.1458980337503154\paperheight,
 bottom = 0.1458980337503154\paperheight
 }
\setlength{\jot}{0.013155617496424828\paperheight}
\linespread{1.1458980337503154}

\setlength{\parskip}{0.013155617496424828\paperheight} % paragraph spacing


%%% COLORS
%%%%%%%%%%%%%%%%%%%%%%%%%%%

\definecolor{red1}{HTML}{f38181}
\definecolor{yellow1}{HTML}{fce38a}
\definecolor{green1}{HTML}{95e1d3}
\definecolor{blue1}{HTML}{66bfbf}
\definecolor{hsblue}{HTML}{00b1db}
\definecolor{hsgrey}{HTML}{afafaf}

%%% CONSTANTS
%%%%%%%%%%%%%%%%%%%%%%%%%%%
\newlength{\smallvert}
\setlength{\smallvert}{0.0131556\paperheight}


%%% COMMANDS
%%%%%%%%%%%%%%%%%%%%%%%%%%%

% differential d
\newcommand*\dif{\mathop{}\!\mathrm{d}}

% horizontal line
\newcommand{\holine}[1]{
  	\begin{center}
	  	\noindent{\color{hsgrey}\hdashrule[0ex]{#1}{1pt}{3mm}}\\%[0.0131556\paperheight]
  	\end{center}
}

% mini section
\newcommand{\minisec}[1]{ \noindent\underline{\textit {#1} } \\}

% quick function plot
\newcommand{\plotfun}[3]{
  \vspace{0.021286\paperheight}
  \begin{center}
    \begin{tikzpicture}
      \begin{axis}[
        axis x line=center,
        axis y line=center,
        ]
        \addplot[draw=red1][domain=#2:#3]{#1};
      \end{axis}
    \end{tikzpicture}
  \end{center}
}

% box for notes
\newcommand{\notebox}[1]{

\tcbset{colback=white,colframe=green1!100!black,title=Note!,width=0.618\paperwidth,arc=0pt}

 \begin{center}
  \begin{tcolorbox}[]
   #1 
  \end{tcolorbox}
 
 \end{center} 
 
}

% box for equation
\newcommand{\eqbox}[2]{
	
	\tcbset{colback=white,colframe=green1!100!black,title=,width=#2,arc=0pt}
	
	\begin{center}
		\begin{tcolorbox}[ams align*]
				#1
		\end{tcolorbox}
		
	\end{center} 
	
}
% END OF PREAMBLE

%%%%%%%%%%%%%%%%%%%%%%%%%%%%%%%%%%%%%

\begin{document}

%%%%%%%%%%%%%%%%%%%%%%%%%%%%%%%%%%%%%
  \includepdf{./titlepage/titlepage.pdf}
  \clearpage
  \setcounter{page}{1}
%%%%%%%%%%%%%%%%%%%%%%%%%%%%%%%%%%%%%

  \section{Wie ist ein C-Programm aufgebaut?}

  \subsection{Header.h}
Mittels Präprozessoranweisung (\inlinecode{#include}) können \emph{Header-Dateien} in
C-Programme eingebunden werden. Diese enthalten \inlinecode{#define} Anweisungen zur
Definition von Konstanten und Macros sowie Funktions- und Strukturprototypen
(Deklarationen) und erlauben somit eine modulare und damit übersichtliche
Strukturierung des Programmes. Außerdem können sie als einfache Dokumentation
dienen, da z.B. Ein- und Ausgaben einer Funktion direkt aus der Deklaration
erkennbar sind.

\subsection{Quellcode.c}
Quellcodedateien enthalten die eigentliche Programmlogik/den Programmablauf. Hierzu gehört die Datei \emph{main.c} sowie alle nötigen \emph{Funktionsdefinitionen}.

Innerhalb der \emph{main.c} werden erst alle Präprozessoranweisungen wie
\inlinecode{#include} oder \inlinecode{#define} gelistet. Als Einstiegspunkt des Programmes gilt dann
die \inlinecode{main()} Funktion, welche dem Betriebssystem oft den Integer-Wert $0$
zur Rückmeldung eines Fehlerfreien Ablaufes zurückgibt; Für Microcontroller ist
dies jedoch nicht von Bedeutung, da sie zum einen kein Betriebssystem besitzen
und zum anderen meist in einer Endlosschleife innerhalb der
\inlinecode{main()}Funktion laufen.


\vspace{\parskip}
\begin{lstlisting}
  #include <stdio.h>
  #define KONST 0
  
  int main() {
   ... 
    return 0;
  }
\end{lstlisting}

  \section{Was sind Assembler, Linker, Compiler?}

  \subsection{Compiler}
  Der Compiler liest den Quellcode, prüft ihn auf syntaktische Fehler und
  übersetzt ihn in Assemblersprache. Die einzelnen
  Quellcodedateien werden in sogenannte \emph{Objektdateien} übersetzt.

  \subsection{Linker}
  Nach dem Kompilieren fügt der Linker die einzelnen Objektdateien sowie
  dazugehörige Bibliotheken zu einer auf
  dem Zielsystem ausführbaren Datei zusammen.

  \subsection{Assembler}
  Der Assembler bildet die Schnittstelle zwischen Assembler- und
  Maschinensprache. Assemblersprache besteht aus symbolischen Abkürzungen,
  welche jeweils stellvertretend für bestimmte Maschinenbefehle stehen.
  Maschinen- und damit Assemblerbefehle sowie deren Syntax sind vom System/der Architektur abhängig und daher nicht
  einheitlich (z.B. Unterschied RISC/CISC)

\vspace{\parskip}
\begin{lstlisting}[language=Assembler]
 MOV R8, #0123 ;lade (MOVe) den Wert #0123 in R8
\end{lstlisting}

  \section{Was ist eine Entwicklungsumgebung?}
  Eine Entwicklungsumgebung ist ein Programm, das mehrere Programme zur
  Softwareentwicklung, wie z.B. einen Programmcode-Editor, Syntax-Highlighting,
  Compiler, Linker, Debugger, GUI, etc., für eine oder mehrere Programmiersprachen bereitstellt.

  \section{Welche Aufgabe haben die Module eines C-Programms?}
  siehe 1: Header + C-dateien
  
\end{document}

% Local Variables:
% TeX-engine: luatex
% End:
